\documentclass{wmecai}
\usepackage[utf8]{inputenc}
\usepackage[brazil]{babel}
\usepackage{amsmath, amssymb}
\usepackage{graphicx}
\usepackage{booktabs}
% Forçar notas de rodapé a serem numeradas
\makeatletter
\renewcommand\@fnsymbol[1]{\ifcase#1\or 1\or 2\or 3\fi}
\makeatother


\title{Otimização da Destinação da Produção de Laranja entre Mercado Interno e Indústria de Suco Considerando Sazonalidade de Preços e Janelas Operacionais de Colheita Utilizando Programação Linear e GurobiPy}
\author{
Alexander Ferreira\thanks{ferreira.alexander@aol.com}, 
Leandro de Andrade\thanks{andradeleandro@usp.br}, 
Afonso Paiva Neto\thanks{apneto@icmc.usp.br}\\
ICMC-USP
}

\date{}

\begin{document}

\criartitulo

\section{Introdução}
O Brasil é o maior produtor mundial de laranja, sendo o estado de São Paulo responsável pela maior parcela da produção nacional. A comercialização da safra representa uma das decisões mais importantes para o produtor rural. Este trabalho aplica técnicas de Pesquisa Operacional para apoiar a decisão de destinação da produção entre mercado interno e indústria de suco.
Este trabalho apresenta um modelo de Programação Linear para otimizar a destinação da produção de laranja entre mercado interno e indústria de suco. A formulação considera a sazonalidade dos preços e as janelas operacionais de colheita, respeitando a perecibilidade da fruta. O modelo foi implementado em Python utilizando GurobiPy e validado por meio de cenários simulados. Os resultados demonstram que o preço pago pela indústria é o fator mais relevante e que a flexibilidade operacional da colheita possui valor econômico mensurável, permitindo ao produtor maximizar sua receita e transformar decisões empíricas em quantitativas.

\section{Problema de Pesquisa}
Um produtor da região de Ribeirão Preto possui produção prevista de 100.000 toneladas de laranja no segundo semestre de 2025. A produção pode ser destinada ao mercado interno (fruta in natura) ou à indústria de suco. A colheita pode ser postergada até 10 dias para o mercado interno e até 20 dias para a indústria. O objetivo é maximizar a receita total da propriedade.

\section{Formulação Matemática}
\subsection{Função Objetivo}
\begin{equation}
Z = \sum_{t \in T} P_{M_t} M_t + \sum_{t \in T} P_{I_t} I_t
\end{equation}
onde $M_t$ é a quantidade destinada ao mercado interno e $I_t$ à indústria no mês $t$.

\subsection{Restrições}
\begin{itemize}
    \item Balanço da produção mensal: $M_t + I_t + T_{M_t} + T_{I_t} = Q_t + T_{M_{t-1}} + T_{I_{t-1}}$
    \item Postergação máxima: $T_{M_t} \leq 0.20 Q_t$, $T_{I_t} \leq 0.40 Q_t$
    \item Capacidade do mercado interno: $M_t \leq 12.000$
    \item Compromisso mínimo com a indústria: $\sum I_t \geq 40.000$
    \item Não negatividade: $M_t, I_t, T_{M_t}, T_{I_t} \geq 0$
\end{itemize}

\section{Implementação Computacional}
O modelo foi implementado em Python utilizando GurobiPy, com apoio das bibliotecas \texttt{pandas} e \texttt{matplotlib}. O método de solução empregado foi o algoritmo Simplex.

\section{Resultados}
Foram simulados quatro cenários:

\begin{table}[h]
\centering
\caption{Resultados dos cenários simulados}
\begin{tabular}{lcccc}
\toprule
Cenário & Receita Total (R\$ milhões) & Mercado (ton) & Indústria (ton) & Receita (\%) \\
\midrule
Base             & 192,24 & 60.000 & 40.000 & 0,00 \\
Mercado -15\%    & 173,25 & 48.000 & 52.000 & -9,87 \\
Indústria +20\%  & 206,09 & 24.000 & 76.000 & 7,20 \\
Janela Reduzida  & 189,65 & 60.000 & 40.000 & -1,34 \\
\bottomrule
\end{tabular}
\end{table}

\noindent As simulações evidenciam dinâmicas cruciais:

\noindent A flexibilidade da colheita tem valor econômico mensurável.  
No Cenário 4 (Janela Reduzida), a receita cai para R\$ 189,65 milhões, uma perda de R\$ 2,59 milhões (-1,34\%) em relação ao Base. Isso mostra que reduzir as janelas de postergação limita a capacidade de capturar preços melhores.

\noindent A rentabilidade é altamente sensível ao preço industrial. 
No Cenário 3 (Indústria +20\%), a receita sobe para R\$ 206,09 milhões, um ganho de R\$ 13,85 milhões (+7,20\%) sobre o Base. O volume destinado à indústria aumenta para 76.000 t, evidenciando que esse canal se torna prioritário quando valorizado.

\noindent O mercado interno é vulnerável a quedas de preço.  
No Cenário 2 (Mercado -15\%), a receita despenca para R\$ 173,25 milhões, uma redução de R\$ 19 milhões (-9,87\%) em relação ao Base. O volume destinado ao mercado cai para 48.000 t, enquanto a indústria absorve 52.000 t.

\noindent O modelo otimizado supera a estratégia tradicional.  
A estratégia empírica de colher e vender imediatamente gera R\$ 197,75 milhões, mas não respeita restrições contratuais e operacionais. Já o modelo otimizado, mesmo com restrições, identifica alocações mais eficientes e cenários de maior ganho.

\begin{figure}[h]
\centering
\includegraphics[width=0.8\textwidth]{grafico.png}
\caption{Receita máxima por cenário}
\label{fig:receita}
\end{figure}


\noindent Os resultados mostram que o preço pago pela indústria é o fator mais relevante. A flexibilidade operacional da colheita possui valor econômico mensurável, e a redução das janelas de colheita gera perda aproximada de R\$ 13,85 milhões, com um ganho puro de 7,20\% de rentabilidade ao produtor, comparando o cenário Base e o cenário Indústia +20\%). 


\section{Conclusão}

O conjunto dos cenários evidenciou que a receita citrícola é sensível às variações de mercado e às restrições operacionais, e o modelo de Programação Linear proposto, resolvido pelo algoritmo Simplex, mostrou-se capaz de representar de forma realista o processo decisório, com a perecibilidade da fruta e as janelas de postergação no pomar incorporadas de forma explícita, oferecendo suporte quantitativo à estratégia ótima de comercialização.




\section*{Agradecimentos}
Agradeço ao corpo docente do MECAI-USP e às instituições IBGE, Fundecitrus e CEPEA/ESALQ-USP pelo fornecimento dos dados utilizados.

\begin{thebibliography}{00}
\bibitem{boldrini} J. L. Boldrini et al. Álgebra Linear e Aplicações, Harbra, 1984.
\bibitem{cuminato} J. A. Cuminato and V. Ruas. Unification of distance inequalities, Comp. Appl. Math., 2014.
\bibitem{jafelice} R. M. Jafelice et al. Study of the dynamics of HIV, Comp. Appl. Math., 2014.
\bibitem{santos} I. L. D. Santos e G. N. Silva. Controle ótimo em escalas temporais, SBMAC, 2013.
\bibitem{cepea} CEPEA – Centro de Estudos Avançados em Economia Aplicada. 
Indicadores agropecuários: Citros, 2026.
\bibitem{fundecitrus} Fundecitrus. Inventário de árvores de citros do cinturão citrícola de São Paulo e triângulo/sudoeste mineiro. 
Technical report, Fundo de Defesa da Citricultura, Araraquara, SP, 2025. 
Acessado em: 16 set. 2026.
\end{thebibliography}
\end{document}

